\clearpage\SourceLeaf{299}
\begin{center}
\textbf{[GEOMETRISCHER ORT}\par\medskip
\textbf{DER SPITZE DES SPHÄRISCHEN DREIECKS}\par\medskip
\textbf{AUF GEGEBENER BASIS MIT GEGEBENEM INHALT.]}
\end{center}

\medskip

\begin{center}
\textsc{Gauss} an \textsc{Schumacher}. Göttingen, 29. Dezember 1841.
\end{center}

\mbox{.\,.\,.\,.\,.\,.\,.} Kennen Sie eine \textls{gedruckte} Lösung der Aufgabe, den geometrischen Ort der Spitze eines sphärischen Dreiecks mit gegebener Basis und gegebenem Flächeninhalt zu finden? Ich hatte neulich Anlass, danach zu suchen, und war etwas verwundert, dass die analytische Lösung auf einem gewissen Umweg zu einem höchst einfachen Ergebnis führt, dessen synthetischer Beweis sich in ein paar Zeilen geben lässt. \mbox{.\,.\,.\,.\,.\,.\,.}

\medskip

\begin{center}
[\textsc{Schumacher} \emph{an} \textsc{Gauss}.\ \emph{Altona, 3. Januar 1842.}]
\end{center}

\{Eine geometrische Lösung des von Ihnen, mein liebster Freund, erwähnten Problems hat vor nicht allzu langer Zeit \textsc{Steiner} in \textsc{Crelle}s Journal (Bd.\ 2.\ p.\ 45) gegeben, wo er aber seine Vorgänger nicht richtig zitiert. \textsc{Lexell}s Aufsatz steht nicht, wie dort angegeben, in Nova Acta Petrop.\ V.\ Pars 1, sondern in Acta Petrop.\ pro 1781, Pars I. \textsc{Euler} ist später in Nova Acta Tom.\ X.\ p.\ 47 (Variae speculationes etc.) darauf zurückgekommen. \textsc{Legendre} hat ihn
